\chapter{Convertible Arbitrage}\label{ch:s2:convertible-arbitrage}

Buy the convertible, short the stock, and the trade is long volatility and credit and short the borrow. In the first quarter of 2005, investors
withdrew more than 20\% of the capital of convertible arbitrage funds. The funds sold convertibles to pay them, and prices fell below theoretical
values, by a median of 2.7\% in mid-May. The funds lost 7.2\% from January to May, about what a 2.7\% cheapening costs at the typical leverage of three
to one. In 2008 the prime brokers who lent to hedge funds cut the leverage they offered. On this chapter's synthetic market, a delta-hedged book of
thirty convertibles, levered three to one, earns 9.8\% a year on its capital in ordinary times. It loses 11.6\% in a planted three-month episode of
forced selling, of which 12.6 points come from cheapening and 6.9 from credit, and it recovers 22.0\% in the six months after. The build is
\texttt{firm.convarb}.

\section{What a convertible holder owns}

A convertible bond (Book~5, chapter~21) is a corporate bond the holder may exchange for a fixed number of shares. It is a straight bond, plus an option
on the shares, minus whatever the issuer's call removes. The whole is exposed to the issuer's default. Book~5's pricer values it on a grid in the share
price, with a default hazard that rises as the share falls: an equity-to-credit link.

The chapter's bond is five years long, with a 2\% coupon, and converts into 2.5 shares (a conversion price of 40). It is callable after two years if
the share is 30\% above the conversion price. The hazard is 3\% a year at a share price of 40, rising in proportion to $40/S$, with 40\% recovery. At
issue, with the share at 40, the pricer values it at 117.7. Its conversion value is 100, and its delta is 1.80 shares per bond
(\cref{fig:s2:convertible-arbitrage:profile}). At a share price of 80 it trades close to its conversion value (204.8 against 200) with a delta of 2.40.

\begin{definition}[Busted convertible]\label{def:s2:convertible-arbitrage:busted}
A \emph{busted convertible}\index{busted convertible} is a convertible whose share price has fallen far below its conversion price, so that the
conversion option is nearly worthless and the bond trades on its credit; its remaining equity sensitivity comes mostly from the link between the share
price and the issuer's default risk.
\end{definition}

At a share price of 15 the bond is busted. It is worth 81.2, below the straight bond's 87.34 at the issue's credit spread, because the hazard has risen
with the fall. Its delta is still 1.27 shares, an equity elasticity of only 0.23, and almost all of that delta is credit.

\begin{omfigure}
\begin{tikzpicture}
\begin{axis}[width=10cm,height=6cm,xlabel={\small share price},ylabel={\small value per 100 face},tick label style={font=\footnotesize},
  xmin=8,xmax=100,ymin=0,ymax=260,grid=major,grid style={gray!25},table/col sep=comma,
  legend style={font=\scriptsize,at={(0.5,-0.25)},anchor=north},legend columns=3]
\addplot[omThm,thick] table[x=share,y=value]{figdata/strategies-2/08-convertible-arbitrage/profile.csv};
\addplot[omDef,thick,dashed] table[x=share,y=conversion]{figdata/strategies-2/08-convertible-arbitrage/profile.csv};
\addplot[gray,thick,dotted] table[x=share,y=floor]{figdata/strategies-2/08-convertible-arbitrage/profile.csv};
\legend{convertible (model),conversion value,straight bond at the issue's spread}
\end{axis}
\end{tikzpicture}

\omcaption{The chapter's five-year convertible at issue, by share price: model value from Book 5's pricer with an equity-to-credit hazard, conversion
value (2.5 shares), and the straight bond discounted at the issue's credit spread. At a share price of 19 and below the bond is worth less than the straight
bond, as the hazard rises with the fall. Data: \texttt{s2\_convarb.profile}.}\label{fig:s2:convertible-arbitrage:profile}
\end{omfigure}

\section{The hedged trade and its Greeks}

\begin{definition}[Convertible delta hedge]\label{def:s2:convertible-arbitrage:delta}
A \emph{convertible delta hedge}\index{convertible delta hedge} is a short position in the issuer's shares equal to the convertible's delta,
rebalanced as the delta changes, so that the position's value does not move with small share-price changes and its P\&L comes from gamma, carry,
credit and the bond's cheapness.
\end{definition}

Agarwal, Fung, Loon and Naik found that a buy-and-hedge strategy, long convertibles with the equity risk hedged, explains much of convertible arbitrage
funds' returns. Market-wide extremes and the supply of new convertibles matter as well. They read the funds' returns as payment for funding issuers
and passing part of the equity risk on to the equity market.

Mitchell, Pedersen and Pulvino describe the market's structure. Convertible arbitrage funds held up to 75\% of the convertible market. Convertibles
are often issued below their fundamental value and go on trading below it, because they are illiquid. The arbitrageur's main risk is that the bond
cheapens further and has to be sold at the wrong moment.

\texttt{firm.convarb} buys one bond on each of the synthetic market's thirty stocks at issue, on day 1,000, with the share at the conversion price.
The bonds cost 3\% below model value, falling to 1\% over six months. Each position is short its delta, rebalanced daily, at a borrow fee of 0.5\% a
year. The book is levered three to one: capital is a third of the bonds' value, 114.19 per bond at purchase. Each day's P\&L is split into buckets
(\cref{lst:s2:convertible-arbitrage:book}): the bond's move at unchanged credit and the stock hedge (together, gamma), the change in the credit
state, the change in cheapness, coupons and borrow.

\section{Credit and borrow}

\begin{definition}[Credit-hedged convertible]\label{def:s2:convertible-arbitrage:credit}
A \emph{credit-hedged convertible}\index{credit-hedged convertible} is a delta-hedged convertible to which the arbitrageur adds credit protection
(credit default swaps, Book~2, chapter~23, or a short position in the issuer's straight bonds) sized to the bond's sensitivity to the issuer's
credit, so that the position keeps its option value and carry without the default risk.
\end{definition}

Credit protection costs a premium. In the synthetic market it pays back the bond's loss from any rise in the hazard multiplier. Each year it costs the
bond's value at risk from a doubled hazard divided by the protection's annuity. The borrow is the other leg's cost. A short position pays the lender a
fee, and when shares are scarce, the fee rises or the loan is recalled.

\section{2005 and 2008}

The synthetic market plants an episode 500 days after purchase, starting with the first crash. For three months the bonds cheapen from 1\% to 6\%
below model value, the hazard doubles and the borrow fee rises to 5\%. Afterwards the cheapness returns to 1\% over six months. On capital:

\begin{center}
\small
\begin{tabular}{@{}lcc@{}}
\toprule
thirty bonds, three to one & delta-hedged & delta- and credit-hedged\\
\midrule
carry outside the episode and recovery (\%/yr) & 9.8 & 6.0\\
three-month episode (\%) & $-11.6$ & $-5.8$\\
\quad of which cheapness & $-12.6$ & $-12.6$\\
\quad credit, net of any hedge & $-6.9$ & 0.0\\
\quad gamma (bond and stock hedge) & 3.7 & 3.7\\
six-month recovery (\%) & 22.0 & 10.6\\
five years (\%) & 51.9 & 30.1\\
\bottomrule
\end{tabular}
\end{center}

In the episode the shares fell with the crash. The bonds lost 17.0\% of capital at unchanged credit, and the short stock gained 20.7\%: the
convexity paid 3.7 points. Cheapening cost 12.6 points and the doubled hazard 6.9. The annual coupon that fell in the window added 5.3, and the
borrow cost 1.1. The credit hedge paid back its 6.9 points in the episode. It cost 21.8\% of capital over five years, and it gave back its share of
the recovery. The book is long everything at once (\cref{fig:s2:convertible-arbitrage:book}): volatility, credit and liquidity. The worst days for each
are the same days.

\begin{omfigure}
\begin{tikzpicture}
\begin{axis}[width=11cm,height=5.8cm,xlabel={\small years after purchase},ylabel={\small cumulative, \% of capital},tick label style={font=\footnotesize},
  xmin=0,xmax=5,ymin=-20,ymax=60,grid=major,grid style={gray!25},table/col sep=comma,
  legend style={font=\scriptsize,at={(0.5,-0.3)},anchor=north},legend columns=5]
\addplot[black,thick] table[x=year,y=total]{figdata/strategies-2/08-convertible-arbitrage/book.csv};
\addplot[omThm,thick] table[x=year,y=cheapness]{figdata/strategies-2/08-convertible-arbitrage/book.csv};
\addplot[omDef,thick] table[x=year,y=credit]{figdata/strategies-2/08-convertible-arbitrage/book.csv};
\addplot[gray,thick,dashed] table[x=year,y=gamma]{figdata/strategies-2/08-convertible-arbitrage/book.csv};
\addplot[gray!60,thick] table[x=year,y=carry]{figdata/strategies-2/08-convertible-arbitrage/book.csv};
\legend{total,cheapness,credit,gamma,coupons less borrow}
\end{axis}
\end{tikzpicture}

\omcaption{The synthetic delta-hedged convertible book, thirty bonds levered three to one: cumulative P\&L as a share of capital and its buckets over
five years. The episode starts at year 2.0 with the first planted crash. Data: \texttt{s2\_convarb.cumulative}.}\label{fig:s2:convertible-arbitrage:book}
\end{omfigure}

The real episodes had their own mechanics. In 2005 the pressure came from investors: redemptions of more than 20\% of capital in one quarter, 28
specialised funds selling 35\% of their convertibles by the end of the year, and a discount largest around the deadlines for redemption notices. In
2008, Mitchell and Pulvino found, the pressure came from the lenders. The imminent failure of prime brokers suddenly cut the leverage hedge funds
could get, and what had looked like long-term financing became short-term. In both years the arbitrageurs who should have bought the cheap bonds were
the ones selling them.

\section{Strategy files}

\begin{strategyfile}{Delta-hedged convertible}\label{strat:s2:convertible-arbitrage:delta}
\sfield{Who pays you, and why} Issuers, who sell convertibles below fundamental value for quick, cheap capital; arbitrageurs fund them and pass the
equity risk to the stock market.
\sfield{Instruments and venues} Convertible bonds (mostly 144A issues); the issuers' shares, borrowed.
\sfield{Signal} Market price against model value; the new-issue discount.
\sfield{Sizing and execution} Delta-hedged, rebalanced as the delta moves; leverage from prime brokers.
\sfield{Costs} Bond spreads, borrow fees, financing.
\sfield{How it dies} Cheapening in forced selling; leverage withdrawn; borrow recalled.
\sfield{Horizon, capacity, infrastructure} Months to years; a convertible pricer, stock loan, prime broker.
\sfield{Backtest honestly} Model values with realistic credit inputs; episodes like 2005 and 2008; financing that can vanish.
\sfield{Sources} Agarwal, Fung, Loon and Naik (2011); Mitchell, Pedersen and Pulvino (2007): $-7.2\%$ for the funds from January to May 2005; this
chapter: 9.8\% a year on capital, $-11.6\%$ in the episode.
\end{strategyfile}

\begin{strategyfile}{Credit-hedged convertible}\label{strat:s2:convertible-arbitrage:credit}
\sfield{Who pays you, and why} As above, less the credit premium paid away.
\sfield{Instruments and venues} Convertibles, shares, and credit default swaps or straight bonds of the issuer.
\sfield{Signal} Cheapness against the cost of protection.
\sfield{Sizing and execution} Protection sized to the bond's credit sensitivity, resized as the share moves.
\sfield{Costs} Protection premiums; the credit default swap's spread.
\sfield{How it dies} Cheapening is not hedged: it cost 12.6 points in the chapter's episode either way.
\sfield{Horizon, capacity, infrastructure} As above; credit trading lines.
\sfield{Backtest honestly} Charge the protection every day, not only when it pays.
\sfield{Sources} This chapter: carry 6.0\% instead of 9.8\%, the episode $-5.8\%$ instead of $-11.6\%$.
\end{strategyfile}

\begin{strategyfile}{Busted convertible credit trade}\label{strat:s2:convertible-arbitrage:busted}
\sfield{Who pays you, and why} Holders who must sell convertibles that no longer behave like equity, and credit buyers who cannot hold them.
\sfield{Instruments and venues} \omterm{def:s2:convertible-arbitrage:busted}{Busted convertibles}; the issuer's straight debt or credit default swaps.
\sfield{Signal} The convertible's yield against the issuer's straight credit.
\sfield{Sizing and execution} Credit-sized; small equity hedges from the equity-to-credit delta.
\sfield{Costs} Illiquid bonds.
\sfield{How it dies} Default: the bond pays its recovery.
\sfield{Horizon, capacity, infrastructure} Months; credit analysis.
\sfield{Backtest honestly} Recoveries and restructurings as they happened.
\sfield{Sources} No performance figure verified.
\end{strategyfile}

\begin{strategyfile}{New-issue convertible}\label{strat:s2:convertible-arbitrage:newissue}
\sfield{Who pays you, and why} Issuers who pay a discount for speed: convertibles are often issued below fundamental value.
\sfield{Instruments and venues} New 144A issues; the issuers' shares.
\sfield{Signal} Issue price against model value.
\sfield{Sizing and execution} Allocations in the book-build; hedged on day one.
\sfield{Costs} Borrow on newly shorted shares.
\sfield{How it dies} Crowded new issues; discounts that shrink.
\sfield{Horizon, capacity, infrastructure} Weeks; relationships with underwriters.
\sfield{Backtest honestly} Allocations were not guaranteed; assume partial fills.
\sfield{Sources} Mitchell, Pedersen and Pulvino (2007); this chapter's new-issue discount of 3\% is an assumption.
\end{strategyfile}

\section{Tutorial: long everything at once}\label{tut:s2:convertible-arbitrage:tut}

\begin{tutorial}
\textbf{Goal.} Price convertibles with Book~5's pricer, hedge them in delta and credit, and run the book through the planted episode.
\textbf{End state:} the table and the two figures.
\begin{tutsteps}
\item \textbf{The book}.
\omcode{code/firm/convarb/firm_convarb.py}{104}{135}{One bond, short its delta, with credit protection optional, attributed day by day.}\label{lst:s2:convertible-arbitrage:book}
\item \textbf{Returns on capital}.
\omcode{code/strategies-2/08-convertible-arbitrage/python/s2_convarb.py}{46}{60}{Carry, the episode by bucket, the recovery.}\label{lst:s2:convertible-arbitrage:summary}
\item \textbf{Run} \texttt{summary()}, \texttt{summary(True)}, \texttt{profile()} and \texttt{fig\_convarb.py}.
\end{tutsteps}
\textbf{What to change next.} Make the borrow recall force a partial unwind of the hedge; lever to five to one and find the episode loss that
triggers a margin call; start the book at a share price of 20 and see what a \omterm{def:s2:convertible-arbitrage:busted}{busted convertible}'s hedge is made of.
\end{tutorial}

\section{Build: convertible arbitrage}\label{bld:s2:convertible-arbitrage:convarb}

\begin{build}
\bfield{Purpose} A convertible book on Book~5's pricer with delta and credit hedges, borrow, and a forced-selling scenario.
\bfield{Interface} \texttt{ArbConfig(\dots)}, \texttt{value\_tables(cfg)}, \texttt{mark(tables, tau, S, stressed)}, \texttt{delta(tables, tau, S,
stressed)}, \texttt{scenario(T, cfg)}, \texttt{run\_book(S, cfg, credit\_hedge)}.
\bfield{Rules} Values tabulated monthly in maturity for a normal and a stressed hazard; market price is model value less cheapness; buckets add up
to the total.
\bfield{Acceptance tests} \texttt{code/firm/convarb/tests/}: the scenario's shape; marks equal the pricer at table nodes, stress lowers the value,
the delta is below the conversion ratio; buckets add up on a calm path with no credit P\&L.
\bfield{Stretch} Borrow recalls; margin calls; several issuers' credit correlated.
\end{build}

\begin{omsources}
\item M.~Mitchell, L.~H.~Pedersen and T.~Pulvino, ``Slow moving capital'', \emph{American Economic Review} 97(2), 2007 (NBER working paper 12877).
\item M.~Mitchell and T.~Pulvino, ``Arbitrage crashes and the speed of capital'', \emph{Journal of Financial Economics} 104(3), 2012.
\item V.~Agarwal, W.~H.~Fung, Y.~C.~Loon and N.~Y.~Naik, ``Risk and return in convertible arbitrage: evidence from the convertible bond market'',
\emph{Journal of Empirical Finance} 18(2), 2011.
\end{omsources}

\section{Exercises}

\begin{exercise}[$\star$]\label{exo:s2:convertible-arbitrage:1}
A bond converts into 2.5 shares. What is its conversion price per 100 of face, and its conversion value at a share price of 60?
\end{exercise}

\begin{exercise}[$\star$]\label{exo:s2:convertible-arbitrage:2}
Bonds cheapen by 2.7\%. What does a fund levered three to one lose on its capital?
\end{exercise}

\begin{exercise}[$\star$]\label{exo:s2:convertible-arbitrage:3}
A bond worth 81.2 has a delta of 1.27 shares at a share price of 15. What is its equity elasticity?
\end{exercise}

\begin{exercise}[$\star\star$]\label{exo:s2:convertible-arbitrage:4}
Why can a \omterm{def:s2:convertible-arbitrage:busted}{busted convertible} be worth less than the straight bond at the issue's credit spread?
\end{exercise}

\begin{exercise}[$\star\star$]\label{exo:s2:convertible-arbitrage:5}
In the episode, the bonds lost 17.0\% of capital and the stock hedge gained 20.7\%. Why did the hedge gain more than the bonds lost?
\end{exercise}

\begin{exercise}[$\star\star$]\label{exo:s2:convertible-arbitrage:6}
Why is cheapening the risk that no hedge removes?
\end{exercise}

\begin{exercise}[$\star\star\star$]\label{exo:s2:convertible-arbitrage:7}
\emph{Coding.} Compare \texttt{summary()} with \texttt{summary(True)}. What does the credit hedge cost over five years, and what does it buy?
\end{exercise}

\begin{exercise}[$\star\star\star$]\label{exo:s2:convertible-arbitrage:8}
\emph{Find the flaw.} ``Our convertible book is delta-hedged and credit-hedged, so it has no risk and we can lever it ten to one.''
\end{exercise}

\section{Problem: Long Everything at Once}

\begin{problem}[{Weekend problem --- a convertible arbitrage book}]\label{pb:s2:convertible-arbitrage:1}
The chapter's synthetic book, Book~5's pricer and the public record.

\medskip\noindent\textbf{Part I --- The bond.}
\begin{enumerate}
\item What does a convertible holder own?
\item Give the chapter's bond's value, conversion value and delta at issue.
\item Define a \omterm{def:s2:convertible-arbitrage:busted}{busted convertible} and describe the bond at a share price of 15.
\item Why does the equity-to-credit link matter for the hedge?
\end{enumerate}

\medskip\noindent\textbf{Part II --- The trade.}
\begin{enumerate}[resume]
\item Define the \omterm{def:s2:convertible-arbitrage:delta}{convertible delta hedge} and the \omterm{def:s2:convertible-arbitrage:credit}{credit-hedged convertible}.
\item What did Agarwal and co-authors find?
\item What did Mitchell, Pedersen and Pulvino find about 2005?
\item What changed in 2008?
\end{enumerate}

\medskip\noindent\textbf{Part III --- The book.}
\begin{enumerate}[resume]
\item How is the synthetic book built, and how is its P\&L attributed?
\item Give its carry, with and without the credit hedge.
\item Decompose the episode's loss.
\item What happened in the recovery?
\end{enumerate}

\medskip\noindent\textbf{Part IV --- The verdict.}
\begin{enumerate}[resume]
\item State the \emph{named result}: the hedged book's carry and its loss in the planted forced-selling episode, and the share from credit.
\item What does the credit hedge cost and buy?
\item Why is leverage the book's real risk?
\item Why were arbitrageurs sellers when bonds were cheapest?
\item How would you size the book?
\item How would you backtest it honestly?
\item Which strategy file depends most on the new-issue discount?
\item In one sentence: what is a convertible arbitrageur paid for?
\end{enumerate}
\end{problem}

\section{Interview questions}

\begin{interviewq}[$\star$ \iqroles{trader}]\label{iq:s2:convertible-arbitrage:1}
Walk through a convertible arbitrage trade.
\end{interviewq}

\begin{interviewq}[$\star\star$ \iqroles{researcher}]\label{iq:s2:convertible-arbitrage:2}
How does the link between the share price and credit change a convertible's delta?
\end{interviewq}

\begin{interviewq}[$\star\star$ \iqroles{risk}]\label{iq:s2:convertible-arbitrage:3}
A convertible book is hedged in delta and credit. What risks remain, and how would you stress them?
\end{interviewq}

\begin{interviewq}[$\star\star$ \iqroles{trader}]\label{iq:s2:convertible-arbitrage:4}
The borrow on one of your short hedges is recalled. What do you do?
\end{interviewq}

\begin{interviewq}[$\star\star$ \iqroles{developer}]\label{iq:s2:convertible-arbitrage:5}
Design the daily process that marks and hedges a book of 200 convertibles.
\end{interviewq}

\begin{interviewq}[$\star\star\star$ \iqroles{researcher}]\label{iq:s2:convertible-arbitrage:6}
Show that a delta-hedged convertible earns about $\tfrac12\Gamma S^2(\sigma_r^2 - \sigma_i^2)\,\mathrm{d}t$ plus carry, and explain what a
default does to the hedge.
\end{interviewq}
